\documentclass[11pt]{article}

\usepackage[margin=1in]{geometry}
\usepackage{amsmath}
\usepackage{amssymb}
\usepackage{booktabs}
\usepackage{enumitem}
\usepackage{hyperref}

\title{KaonLT Particle-Background Subtraction Procedure}
\author{KaonLT Production Analysis}
\date{\today}

\begin{document}
\maketitle

\section{Purpose and Scope}

This note documents the current particle-background treatment in the KaonLT
production analysis.  The nominal kaon path uses two separate corrections:
\begin{enumerate}[leftmargin=2em]
  \item event-level slow-proton cleaning from coincidence timing, SHMS
        $\delta$, and aerogel response; and
  \item an MM-dependent pion-background subtraction from aligned SIMC neutron,
        $\Delta$, and SIDIS component shapes.
\end{enumerate}

The pion fit uses the kaon spectrum after slow-proton cleaning.  The standard
processing order is
\begin{align}
  \text{random/dummy correction}
  \longrightarrow \text{freeze bin-counting spectra}
  \longrightarrow \text{slow-proton cleaning}
  \longrightarrow \text{pion component subtraction}
  \longrightarrow \text{optional empirical residual fits}.
\end{align}

The main code owners are
\begin{itemize}[leftmargin=2em]
  \item \texttt{src/cuts/rand\_sub.py} for setting-wide ordering and
        application;
  \item \texttt{src/cuts/proton\_contamination\_weights.py} for the
        slow-proton method;
  \item \texttt{src/cuts/pion\_component\_fits.py} and
        \texttt{src/cuts/pion\_component\_subtraction.py} for pion templates
        and weights;
  \item \texttt{src/binning/calculate\_yield.py} and
        \texttt{src/binning/ave\_per\_bin.py} for bin-level propagation; and
  \item \texttt{src/utility/background\_config.py} for defaults and
        setting/$\phi$ overrides.
\end{itemize}

\section{Common Data Staging}

For each $\phi$ setting, prompt, random, dummy, and dummy-random sources are
combined with their usual signs and normalization coefficients.  For any
observable $X$,
\begin{align}
  H_X^{\mathrm{data}}
  ={}&c_{\mathrm{prompt}}H_X^{\mathrm{prompt}}
     +c_{\mathrm{rand}}H_X^{\mathrm{rand}} \\
    &+c_{\mathrm{dummy}}H_X^{\mathrm{dummy}}
     +c_{\mathrm{dummy\ rand}}H_X^{\mathrm{dummy\ rand}},
\end{align}
where the random and dummy terms carry the appropriate negative signs.

Before either particle correction, the kaon $t$ and $\phi$ distributions used
for bin finding are cloned.  The binning definition therefore remains fixed,
while the final histograms, yields, and averages include both corrections.

\section{Part I: Slow-Proton Cleaning}

\subsection{Method and Input Sample}

Slow protons can survive the kaon selection as a displaced timing component.
They are removed with the configured
\texttt{ctime\_aero\_event\_weight} method, implemented as
\texttt{c\_script\_exact}.  The Python implementation follows the accepted
\texttt{check\_slow\_protons.C} procedure; it is not a globally scaled
proton-template subtraction.

The fit uses explicit no-RF kaon trees for prompt, random, dummy, and
dummy-random events.  Their signed, prompt-relative count-scale combination is
used to build the timing-fit sample before pion subtraction.  Corrected physics
targets are subsequently filled with the physical source coefficients.

The timing sample requires the standard kaon acceptance, HGCer-hole rejection,
and the no-MM-cut selection.  In strict mode, a missing explicit no-RF input is
an error rather than a reason to substitute an RF-selected tree.  The stored
result records the resolved setting, $\phi$ setting, override layers, input
tree policy, and failure policy.

\subsection{Timing Fits}

The default PID partition is
\begin{align}
  N_{\mathrm{aero}}&\in[0,3),[3,6),[6,10),[10,15),[15,25), \\
  \delta_{\mathrm{SHMS}}&\in[-10,20]\ \% \quad\text{in ten bins}.
\end{align}
For each aerogel slice, a global timing fit determines the kaon and proton
shapes.  Per-$\delta$ fits then use those shapes as constraints while finding
the local component normalizations and timing behavior.  The model contains
kaon and proton Gaussian components and, when needed, a non-negative
residual-other component.

Fit quality, peak separation, component significance, support, and bound-hit
tests determine whether an aerogel--$\delta$ cell is usable.  The default
timing policy, \texttt{rf\_then\_ct\_best}, tests usable RF-like branches and
\texttt{CTime\_ROC1}; it retains the preferred valid result.  The diagnostics
record any per-aerogel multi-start, standard-fit, pre-diamond, or local-peak
rescue attempts used to make that decision.

\subsection{Event Weights and RF Handling}

For timing value $\tau$, the proton probability is
\begin{align}
  P_p(\tau,\delta,N_{\mathrm{aero}})
  =\frac{f_p(\tau)}
         {f_p(\tau)+f_K(\tau)+f_{\mathrm{other}}},
  \qquad 0\leq P_p\leq1.
\end{align}
For a source with physical coefficient $c_s$, the corresponding proton and
cleaned contributions are
\begin{align}
  w_{p,s}=c_sP_p,\qquad
  w_{\mathrm{clean},s}=c_s(1-P_p).
\end{align}
The output is therefore a parallel set of raw, estimated-proton, and
proton-cleaned histograms.  The per-event factors are frozen in a lookup keyed
by source and entry number, then reused in the setting-wide, yield, and
averaging paths.

For low-$\epsilon$ data, RF membership is optionally restored only after the
proton factor is evaluated, by signature matching to the existing RF skim
trees.  The default policy is \texttt{epsset\_default\_after\_cleaning};
high-$\epsilon$ data do not receive this extra RF-membership step.  Thus RF
selection cannot change the no-RF sample used to determine $P_p$.

\subsection{Acceptance and Diagnostics}

Each $\delta$ bin is classified as supported, marginal, or unsupported from
its usable aerogel cells and modeled yield.  At least one supported or marginal
bin is required for an accepted cleaning result.  An unaccepted result leaves
the particle targets unchanged, while a missing required no-RF or RF-reference
input remains a strict-mode error where configured.

The diagnostic PDF and serialized result contain
\begin{itemize}[leftmargin=2em]
  \item selected timing branch, global timing shapes, and fit constraints;
  \item aerogel and $\delta$ support, fitted yields, and proton fractions;
  \item event-probability closure: summed event proton probability versus
        fitted proton yield;
  \item average proton weights versus $\delta$ and
        $(\delta,N_{\mathrm{aero}})$; and
  \item raw, proton-estimated, cleaned, and final-RF MM spectra, including
        low-MM and $\Lambda$-peak validation windows.
\end{itemize}

\section{Part II: Pion SIMC-Shape Subtraction}

\subsection{Templates and Provenance}

After slow-proton cleaning, the nominal particle-subtraction mode is
\texttt{simc\_shape\_components}.  The components are
\begin{itemize}[leftmargin=2em]
  \item $\pi$-$n$,
  \item $\pi$-$\Delta$, and
  \item $\pi$-SIDIS.
\end{itemize}

Their ROOT inputs are resolved through the runtime
\texttt{background\_samples} payload and the setting/$\phi$ map, not
hard-coded in the fit.  The normal SIMC tree is \texttt{h10}.  Each component
is filled with SIMC event weights and comparison cuts, then normalized to a
unit-area MM template.

Templates can be resolved at setting-wide and $(t,\phi)$ scope.  For a sparse
fine bin, the fallback sequence is direct $(t,\phi)$, then $t$-aggregated, then
setting-wide.  This is a template-support fallback and does not replace the
configured physical input sample.

\subsection{Dynamic Template Alignment}

Pion-control data calibrate the SIMC MM alignment before amplitudes are fit.
This dynamic-alignment layer is independent of the legacy residual-shift
switches.  At setting-wide scope, $\pi$-$n$ scans offsets from $-0.020$ to
$+0.020$~GeV in $0.001$~GeV steps, while $\pi$-$\Delta$ scans $-0.030$ to
$+0.030$~GeV in $0.002$~GeV steps.  Candidate fit-window expansions are also
tested.

The selected alignment must satisfy support, retained-template-integral,
score-improvement, and localized-minimum requirements.  Offset-boundary hits
are rejected by default; expansion-boundary hits are diagnostic warnings by
default.  The accepted pion-control map is shared with the kaon fit.  Fine
bins can scan around their parent map, and only rejected components inherit
the parent or common fallback.  $\pi$-SIDIS currently inherits the parent map
by default rather than making an independent scan.

The saved payload keeps proposed scan metrics distinct from the metrics of the
alignment actually applied.  This distinction is important when a mixed
fine-bin map uses an accepted component-specific scan together with parent
fallback for another component.

\subsection{Fit Targets and Current Configuration}

The aligned templates are fit to
\begin{itemize}[leftmargin=2em]
  \item pion-control data, yielding denominator amplitudes
        $B_n$, $B_\Delta$, and $B_{\mathrm{SIDIS}}$; and
  \item proton-cleaned kaon no-subtraction data, yielding numerator amplitudes
        $A_n$, $A_\Delta$, and $A_{\mathrm{SIDIS}}$.
\end{itemize}

Both paths use non-negative staged amplitudes: each component is fit to the
residual after previous components.  Table~\ref{tab:pion-config} is the
current nominal snapshot.  The resolved configuration, including a
setting/$\phi$ override, remains authoritative for a production result.

\begin{table}[htbp]
  \centering
  \caption{Nominal active pion-component fit configuration.}
  \label{tab:pion-config}
  \begin{tabular}{@{}lll@{}}
    \toprule
    Target & Fit mode and order & Active fit windows (GeV) \\
    \midrule
    Pion control &
      staged plus regularized joint: $\pi n\to\pi$-SIDIS$\to\pi\Delta$ &
      $\pi n:[0.88,0.94]$; $\pi\Delta:[1.16,1.18]\cup[1.24,1.26]$;
      $\pi$-SIDIS$:[1.05,1.10]$ \\
    Kaon no-sub &
      staged only: $\pi n\to\pi$-SIDIS$\to\pi\Delta$ &
      $\pi n:[0.88,0.94]$; $\pi\Delta:[1.16,1.18]\cup[1.24,1.26]$;
      $\pi$-SIDIS$:[1.05,1.10]$ \\
    \bottomrule
  \end{tabular}
\end{table}

The $K\Sigma^0$ stage and its $[1.17,1.23]$~GeV window remain configured but
are disabled in both current fits.  The pion-control path may regularize a
joint non-negative refinement around its staged solution.  The kaon
no-subtraction path is presently staged-only.  All active post-fit and
post-refinement component scales are unity, so the control-model denominator
is not phenomenologically rescaled.

For a least-squares staged component $T_j$ with residual $R_{j-1}$,
\begin{align}
  C_j=\max\left(0,
    \frac{\sum_iT_j(i)R_{j-1}(i)/\sigma_i^2}
         {\sum_iT_j(i)^2/\sigma_i^2}\right),
  \qquad R_j=R_{j-1}-C_jT_j.
\end{align}
Components configured for window-integral normalization instead use the ratio
of residual and template integrals in their amplitude window.  A joint
refinement augments the data objective with priors around the staged result
and retains non-negative amplitudes.

\subsection{MM-Dependent Event Weight}

Let the aligned unit-normalized templates be
\begin{align}
  T_n(MM),\qquad T_\Delta(MM),\qquad T_{\mathrm{SIDIS}}(MM).
\end{align}
The fitted control and kaon-side models are
\begin{align}
  P_{\mathrm{ctrl}}^{\mathrm{fit}}(MM)
  &=B_nT_n(MM)+B_\Delta T_\Delta(MM)
    +B_{\mathrm{SIDIS}}T_{\mathrm{SIDIS}}(MM),\\
  K_{\pi\mathrm{-bg}}^{\mathrm{fit}}(MM)
  &=A_nT_n(MM)+A_\Delta T_\Delta(MM)
    +A_{\mathrm{SIDIS}}T_{\mathrm{SIDIS}}(MM).
\end{align}
The pion event weight is
\begin{align}
  w_\pi(MM)
  =\frac{K_{\pi\mathrm{-bg}}^{\mathrm{fit}}(MM)}
         {P_{\mathrm{ctrl}}^{\mathrm{fit}}(MM)}.
\end{align}
It is evaluated with the configured denominator floor and clip policy.  Its
model-level closure identity is
\begin{align}
  w_\pi(MM)P_{\mathrm{ctrl}}^{\mathrm{fit}}(MM)
  =K_{\pi\mathrm{-bg}}^{\mathrm{fit}}(MM).
\end{align}

The physical subtraction object is the weighted pion-control event template:
$w_\pi$ is applied to accepted pion-control prompt, random, dummy, and
dummy-random events with their usual signs and normalization coefficients.
The final kaon target is therefore
\begin{align}
  H_{\mathrm{kaon,after\ pion}}
  =H_{\mathrm{kaon,proton\ cleaned}}
   -H_{\mathrm{weighted\ pion\ event\ template}}.
\end{align}
The weighted pion-control model is an algebra check; the weighted event
template is the object actually subtracted from data.

\subsection{Fallback and Diagnostics}

The nominal production fallback for component subtraction is \texttt{error}.
An invalid component result must not silently become a legacy single-scale
subtraction.  The low-$\epsilon$ optimizer prepass is an intentional special
case that may request \texttt{single\_scale} so a sparse exploratory bin does
not abort the full scan; it does not redefine the nominal production mode.

Pion diagnostics distinguish alignment scans and applied maps, staged and
joint-fit models, model closure, event-template closure, and kaon MM spectra
before and after pion subtraction.  Event-template closure is the physics
validation: it checks whether the weighted control events reproduce the
kaon-side pion background that is actually being removed.

\section{Reproducibility Record}

Each setting-level result should retain the resolved particle-subtraction
configuration, background-sample provenance, dynamic pion-alignment payload,
serialized slow-proton result, accepted or fallback status, and diagnostic
PDF pages.  These records trace a final yield through
\begin{align}
  \text{raw kaon data}
  \to\text{random/dummy corrected data}
  \to\text{proton-cleaned data}
  \to\text{pion-subtracted data}.
\end{align}

\section{Summary}

The current analysis first applies an event-level timing--$\delta$--aerogel
slow-proton correction, then derives an MM-dependent pion weight from
dynamically aligned SIMC component fits.  Both corrections are propagated to
the same event sources and final targets, while bin-counting inputs are fixed
before either correction.  This keeps the two physical backgrounds, their
validations, and their configuration ownership explicit.

\end{document}
